% !TEX root = ../main.tex
\corrigesChapitre{9}{Synthèses organiques}
\corrigesSouples

\begin{blocLivre}{cor-9-1}
\begin{corrige}{1}{Les étapes d'un protocole}
\demarche{Transformer, isoler le produit brut, le purifier, puis l'analyser.}
\begin{questions}
\item Transformation : \res{\text{b}} ; isolement : \res{\text{f, c}} ; purification : \res{\text{d, g}} ; analyse : \res{\text{a, e}}.
\item \res{\text{b, f, c, d, g, a, e}} (les analyses a et e peuvent être échangées).
\item L'aspirine est \res{\text{peu soluble dans l'eau froide}} : elle précipite, et l'on en perd moins lors de la filtration.
\item La \res{\text{température de fusion}} et la \res{\text{CCM}}.
\end{questions}
\end{corrige}
\end{blocLivre}
\begin{blocLivre}{cor-9-2}
\begin{corrige}{2}{Le montage à reflux}
\demarche{Chaque élément a un rôle : contenir, chauffer, condenser, agiter.}
\begin{questions}
\item A : \res{\text{ballon}} (mélange réactionnel) ; B : \res{\text{réfrigérant à eau}} ; C : \res{\text{entrée d'eau}} ; D : \res{\text{sortie d'eau}} ;
E : \res{\text{chauffe-ballon}}.
\item L'eau entre en bas (C) et sort en haut (D) : ainsi le réfrigérant reste \res{\text{entièrement rempli}} d'eau froide.
\item Chauffer \res{\text{accélère}} la transformation ; le reflux \res{\text{évite les pertes}} de réactifs et de produits par évaporation.
\item Le chauffage produit des vapeurs : un montage fermé subirait une \res{\text{surpression}} et pourrait exploser.
\end{questions}
\schema[corrige]{09-ex02}
\end{corrige}
\end{blocLivre}
\begin{blocLivre}{cor-9-3}
\begin{corrige}{3}{Choisir une technique}
\demarche{État physique, solubilité, densité et température d'ébullition guident le choix.}
\begin{questions}
\item \res{\text{Filtration sous pression réduite}} : le solide, peu soluble à froid, reste sur le filtre.
\item \res{\text{Décantation}} dans une ampoule : les deux liquides non miscibles forment deux phases.
\item \res{\text{Séchage}} avec un desséchant (sulfate de magnésium anhydre), puis filtration.
\item \res{\text{Distillation}} : l'impureté, plus volatile, distille la première.
\end{questions}
\end{corrige}
\end{blocLivre}
\begin{blocLivre}{cor-9-4}
\begin{corrige}{4}{Exploiter une chromatographie}
\demarche{$R_\mathrm{f}$ : distance parcourue par la tache divisée par celle de l'éluant, depuis la ligne de dépôt.}
\begin{questions}
\item L'éluant a parcouru \qty{6,0}{cm}. Produit : $\dfrac{3{,}0}{6{,}0}=\res{0{,}50}$ ; autre espèce : $\dfrac{4{,}8}{6{,}0}=\res{0{,}80}$.
\item \res{\text{Non}} : il donne deux taches ; la seconde est à la hauteur du dépôt 3, c'est du \res{\text{réactif}} qui n'a pas réagi.
\item Il ne donne qu'une tache, à la hauteur de la référence : il est \res{\text{pur}} (pour cette technique) et \res{\text{identifié}}.
\item La recristallisation a \res{\text{éliminé le réactif}} resté dans le produit brut.
\end{questions}
\schema[corrige]{09-ex04}
\retiens{Les distances se mesurent toujours depuis la ligne de dépôt : jusqu'au centre de la tache pour l'espèce, jusqu'au
front pour l'éluant. Une espèce pure ne donne qu'une tache ; deux taches à la même hauteur indiquent la même espèce.}
\end{corrige}
\end{blocLivre}
\begin{blocLivre}{cor-9-5}
\begin{corrige}{5}{Un premier rendement}
\demarche{Les nombres stœchiométriques valent 1 : le réactif limitant est le moins abondant.}
\begin{questions}
\item L'\res{\text{alcool}} (\qty{0,100}{mol} contre \qty{0,200}{mol}).
\item $n_\text{max}=\res{\qty{0,100}{mol}}$.
\item $n=\dfrac{9{,}1}{130{,}0}=\qty{0,0700}{mol}$, $\eta=\dfrac{0{,}0700}{0{,}100}=\res{\qty{70}{\percent}}$.
\item La transformation \res{\text{n'est pas totale}} (chapitre~4) ; une partie du produit est \res{\text{perdue}} lors des transferts et de la
purification.
\end{questions}
\retiens{Le rendement se calcule à partir du réactif limitant : c'est lui qui fixe la quantité maximale de produit.}
\end{corrige}
\end{blocLivre}
\begin{blocLivre}{cor-9-6}
\begin{corrige}{6}{Calculer un rendement en Python}
\demarche{\code{m / M / a} est l'avancement qui épuiserait le réactif ; le plus petit est l'avancement maximal.}
\begin{questions}
\item La quantité du réactif divisée par son nombre stœchiométrique : l'avancement qui l'épuiserait. Le \res{\text{minimum}} est
$x_\text{max}$, fixé par le réactif limitant.
\item \code{xmax = 0.01449 mol ; m\_max = 2.61 g} puis \code{rendement = 92.0 \%}.
\item On écrit \code{reactifs = \{'acide': [1, 21.0, 60.0], 'alcool': [1, 12.15, 88.0]\}} et
\code{M\_produit, m\_obtenue = 130.0, 10.9}.\newline
Affichage : \code{xmax = 0.13807 mol ; m\_max = 17.95 g} puis \code{rendement = 60.7 \%}.
\item \code{if m\_obtenue > m\_max: print('erreur : rendement impossible')}.
\end{questions}
\end{corrige}
\end{blocLivre}
\begin{blocLivre}{cor-9-7}
\begin{corrige}{7}{L'aspirine et son banc Kofler}
\demarche{Une impureté abaisse la température de fusion d'un solide.}
\begin{questions}
\item \qty{128}{\degreeCelsius} est inférieur à \qty{135}{\degreeCelsius} : le produit brut est \res{\text{impur}}.
\item On dépose quelques cristaux sur le banc, chaud d'un côté et plus froid de l'autre, et l'on repère la limite entre
solide et liquide. L'étalonnage avec une substance connue \res{\text{corrige}} les écarts de l'appareil.
\item Les impuretés et l'eau \res{\text{alourdissent}} le produit brut : on surestimerait le rendement.
\item $\dfrac{2{,}15}{0{,}500}=4{,}3$ : \res{4 \text{ comprimés}} complets.
\end{questions}
\end{corrige}
\end{blocLivre}
\begin{blocLivre}{cor-9-8}
\begin{corrige}{8}{La recristallisation}
\demarche{Le produit est très soluble à chaud et peu soluble à froid ; les impuretés restent dissoutes.}
\begin{questions}
\item $V=\dfrac{5{,}0}{60}\approx\res{\qty{83}{mL}}$.
\item $3{,}0\times0{,}100=\res{\qty{0,30}{g}}$.
\item $5{,}0-0{,}30=\res{\qty{4,7}{g}}$ au maximum (si le produit brut était uniquement de l'acide benzoïque).
\item Plus il y a de solvant, plus il reste de produit dissous à froid : on en \res{\text{perdrait davantage}}. Les impuretés, très
solubles, \res{\text{restent dans le filtrat}}.
\end{questions}
\retiens{Une recristallisation exploite la différence de solubilité du produit à chaud et à froid : il recristallise en
refroidissant, tandis que les impuretés, peu abondantes ou très solubles, restent dissoutes.}
\end{corrige}
\end{blocLivre}
\begin{blocLivre}{cor-9-9}
\begin{corrige}{9}{Synthétiser l'acide benzoïque}
\demarche{Une mole d'alcool donne une mole d'acide : $n_\text{max}=n_0(\text{alcool})$, l'oxydant étant en excès.}
\begin{questions}
\item Réactif : \res{\text{alcool}} ; produit : \res{\text{acide carboxylique}}.
\item De chaque côté : 7 C, 10 H, 2 O et une charge nulle. L'alcool est le réducteur du couple : il est \res{\text{oxydé}}.
\item $m=2{,}0\times1{,}04=\qty{2,08}{g}$, $n=\qty{1,93e-2}{mol}$ ; $m_\text{max}=\num{1.93e-2}\times122{,}0=\res{\qty{2,35}{g}}$ et $\eta=\dfrac{1{,}80}{2{,}35}\approx\res{\qty{77}{\percent}}$.
\item Par \res{\text{filtration sous pression réduite}} : l'acide benzoïque, peu soluble dans l'eau froide, est solide.
\end{questions}
\end{corrige}
\end{blocLivre}
\begin{blocLivre}{cor-9-10}
\begin{corrige}{10}{Purifier par distillation}
\demarche{L'espèce de plus basse température d'ébullition se vaporise la première.}
\begin{questions}
\item La \res{\text{propanone}}, la plus volatile (\qty{56}{\degreeCelsius}).
\item Il mesure la température des \res{\text{vapeurs}} qui passent dans le réfrigérant, donc identifie l'espèce qui distille.
\item Palier vers \qty{56}{\degreeCelsius} pendant que la propanone distille, montée, palier vers \qty{78}{\degreeCelsius} (éthanol), puis vers
\qty{100}{\degreeCelsius}.
\item Les vapeurs condensées doivent \res{\text{s'écouler vers l'erlenmeyer}}, et non retomber dans le ballon.
\end{questions}
\schema[corrige]{09-ex10-cor}
\end{corrige}
\end{blocLivre}
\begin{blocLivre}{cor-9-11}
\begin{corrige}{11}{Le paracétamol}
\demarche{L'anhydride étant en excès, le 4-aminophénol est limitant : $n_\text{max}=n_0(\text{4-aminophénol})$.}
\begin{questions}
\item De chaque côté : 10 C, 13 H, 1 N et 4 O : l'équation est \res{\text{ajustée}}.
\item $n=\dfrac{3{,}00}{109{,}0}=\qty{2,75e-2}{mol}$ et $m_\text{max}=\num{2.75e-2}\times151{,}0\approx\res{\qty{4,16}{g}}$.
\item $\eta=\dfrac{3{,}20}{4{,}16}\approx\res{\qty{77}{\percent}}$.
\item Elle est égale à celle des tables : le produit est \res{\text{pur}} et probablement du paracétamol. Une \res{\text{CCM}} avec une
référence, ou un spectre infrarouge, confirmerait l'identification.
\end{questions}
\end{corrige}
\end{blocLivre}
\begin{blocLivre}{cor-9-12}
\begin{corrige}{12}{L'arôme de banane}
\demarche{Calculer les deux quantités initiales ; l'alcool, en défaut, fixe $n_\text{max}$.}
\begin{questions}
\item Étape 1 : \res{\text{transformation}} ; étape 2 : \res{\text{isolement}} ; étape 3 : \res{\text{purification}}. L'acide sulfurique est un
\res{\text{catalyseur}}.
\item La phase organique ($\rho=\qty{0,87}{g.mL^{-1}}$) est \res{\text{au-dessus}} de l'eau salée ($\qty{1,2}{g.mL^{-1}}$). L'ester est encore moins
soluble dans l'eau salée que dans l'eau : on en \res{\text{perd moins}}.
\item Il élimine l'\res{\text{acide éthanoïque}} restant, qui réagit en dégageant du \res{\text{dioxyde de carbone}}.
\item $n(\text{alcool})=\dfrac{15{,}0\times0{,}81}{88{,}0}=\qty{0,138}{mol}$ ; $n(\text{acide})=\dfrac{20{,}0\times1{,}05}{60{,}0}=\qty{0,350}{mol}$ : l'alcool est limitant.
$m_\text{max}=0{,}138\times130{,}0=\qty{17,9}{g}$ ; $m=12{,}5\times0{,}87=\qty{10,9}{g}$ ; $\eta\approx\res{\qty{61}{\percent}}$.
\end{questions}
\piege{Calculer le rendement avec le volume de liquide obtenu : il faut d'abord le convertir en masse, $m=\rho\times V$, puis
comparer cette masse à $m_\text{max}$, calculée à partir du réactif limitant.}
\end{corrige}
\end{blocLivre}
\begin{blocLivre}{cor-9-13}
\begin{corrige}{13}{Le savon à froid}
\demarche{Trois moles de soude par mole d'oléine ; le surgras rend la soude limitante.}
\begin{questions}
\item De chaque côté : 57 C, 107 H, 9 O et 3 Na : l'équation est \res{\text{ajustée}}.
\item $n(\text{oléine})=\dfrac{500}{884{,}0}=\qty{0,566}{mol}$, d'où $m(\ce{NaOH})=3\times0{,}566\times40{,}0\approx\res{\qty{67,9}{g}}$ ; avec le surgras :
$0{,}95\times67{,}9\approx\res{\qty{64,5}{g}}$.
\item La \res{\text{soude}} est limitante : elle est entièrement consommée. Il ne reste pas de soude corrosive dans le savon, mais un
peu d'huile, qui \res{\text{nourrit la peau}}.
\item $m(\text{savon})=3\times0{,}566\times304{,}0\approx\res{\qty{516}{g}}$ et $m(\text{glycérol})=0{,}566\times92{,}0\approx\res{\qty{52,0}{g}}$. Avec le surgras, on obtient
\qty{5}{\percent} de savon en moins.
\end{questions}
\end{corrige}
\end{blocLivre}
\begin{blocLivre}{cor-9-14}
\begin{corrige}{14}{Améliorer un rendement}
\demarche{Un catalyseur accélère sans changer l'état final ; un excès de réactif ou l'élimination d'un produit déplace l'état final.}
\begin{questions}
\item $\res{\ce{C2H4O2 + C2H6O -> C4H8O2 + H2O}}$.
\item \res{\text{Non}} : il permet d'atteindre plus vite l'état final, mais le même état final.
\item Le rendement \res{\text{augmente}} avec l'excès d'acide : $\qty{67}{\percent}$, puis $\qty{85}{\percent}$ et $\qty{90}{\percent}$. On met en excès le réactif
\res{\text{le moins coûteux}}, pour transformer au mieux le plus cher.
\item Sans l'eau, la réaction inverse, de l'ester avec l'eau, ne peut plus se produire : la transformation devient
\res{\text{quasi totale}}.
\end{questions}
\end{corrige}
\end{blocLivre}
\begin{blocLivre}{cor-9-15}
\begin{corrige}{15}{Préparer vingt comprimés}
\demarche{On part du produit voulu et l'on remonte : $n_\text{max}=\dfrac{n_\text{obtenue}}{\eta}$, puis $n_0(\text{acide salicylique})=n_\text{max}$.}
\begin{questions}
\item $m=20\times0{,}500=\res{\qty{10,0}{g}}$ et $n=\dfrac{10{,}0}{180{,}0}=\res{\qty{5,56e-2}{mol}}$.
\item $n_\text{max}=\dfrac{\num{5.56e-2}}{0{,}80}=\res{\qty{6,94e-2}{mol}}$.
\item L'acide salicylique est limitant et réagit mole à mole : $m=\num{6.94e-2}\times138{,}0\approx\res{\qty{9,58}{g}}$.
\item Un meilleur rendement consomme \res{\text{moins de réactifs}} et produit moins de déchets pour la même quantité de médicament.
\end{questions}
\retiens{Quand on connaît le rendement, on remonte de la quantité de produit voulue à la quantité de réactif à utiliser :
$n_\text{max}=\dfrac{n_\text{voulue}}{\eta}$.}
\end{corrige}
\end{blocLivre}
